% Two independent successors to 109; together they give 107 instructions.
\section{A positive coefficient and a triangular Pell equation}
\label{sec:final-two-savings}

Two independent changes remove another addition and another multiplication.
The first uses positivity of an already computed coefficient to recover
a bound that previously required an explicit addition. The second
rewrites the first Pell equation using its odd root and proves that
its index modulus can be halved. Each separately gives 108 instructions;
together they give 107.

\subsection{Removing the square from the explicit bound}

Keep the fixed input-unit encoding of Section~\ref{sec:input-unit},
with $\mathcal C=x+g$ and $Y=\ell+eq$. Replace
\[
 Y+\mathcal C^2+\alpha=q^2
\]
by the equations
\begin{equation}\label{eq:positive-coefficient-bound}
 Y+\alpha=q^2,\qquad \Omega=Z\lambda-2e,
 \qquad \Omega>0.
\end{equation}
Here $\Omega$ is a new positive system unknown. Its defining value
is the negative of an existing register. Reverse that subtraction and
compute
\[
 P_1=(Z\lambda-2e)\mathcal C^2,
 \qquad P_2=B\lambda(1+q),\qquad S_3=P_2-P_1.
\]
The value of $S_3$ is unchanged. Reversing these signs has equal
arithmetic cost, and $\Omega$ is connected to the computed coefficient
by a free equality test. Retain the separate positive unknown
$\sigma=S_3$.

\begin{proposition}\label{prop:positive-coefficient108}
The replacements \eqref{eq:positive-coefficient-bound} give a
membership-equivalent system with 108 arithmetic instructions:
60 multiplications and 48 additions. It has 34 positive unknowns
and 22 equations.
\end{proposition}
\begin{proof}
Before decoding, the new bound gives $\ell<Y<q^2$ and $e<q$.
The geometric equation still gives $B\le q^2$ and $q>4b\ge8$.
Consequently
\[
 P_2=B\lambda(1+q)<2q^2(1+q)<3q^3.
\]
Since $\Omega$ and $\sigma$ are positive integers,
\[
 0<S_3=P_2-\Omega\mathcal C^2<P_2,
 \qquad \mathcal C^2<P_2/\Omega\le P_2<3q^3<q^4.
\]
Thus $g<\mathcal C<q^2$. This weaker bound is enough for the first
packing block, whose width is $q^2$. Together with $Y<q^2$ and
$S_3<q^4$, it gives $0<S<n=q^8$.

The mask bounds use only $\ell<q^2$, $B\le q^2$, and
$b<(B-Z)\lambda<q^2-1$. They are unchanged. Explicitly, if
$d_0=\lfloor(b-1)\ell/q^2\rfloor$ and
$\varepsilon_0=(b-1)\ell-d_0q^2$, then
\[
 T=(q^2-1-\varepsilon_0)
       +q^2\bigl((B-Z)\lambda-d_0\bigr)+q^4(B-4)\ell
\]
has normalized blocks in $[0,q^2)$, $[0,q^2)$, and $[0,q^4)$.
Hence $0\le T<n$ and $r\ge n^2-1\ge n$.
The relaxed Pell bootstrap proved in Section~\ref{sec:input-unit}
now applies: it uses $3L\le B\le n\le r$, $q,b\le n$, and the
unchanged Pell equations, without a further bound on $g$ or
$\mathcal C$. It gives $q=B^L$, the required powers of two,
and $\tau_2(S,T)=0$.

We recover the stronger bound before using it in digit decoding.
Now $q=B^L>B$, so $\lambda>q$. Since $e<q$ and $Z\ge4$,
\[
 2e<2q<2\lambda\le Z\lambda/2,
 \qquad \Omega>Z\lambda/2.
\]
It follows that
\[
 \mathcal C^2<\frac{P_2}{\Omega}
       <\frac{2B(1+q)}Z
       \le\frac{B(1+q)}2
       <\frac{q(q+1)}2<q^2.
\]
Thus $g<\mathcal C<q$ holds before the first-mask decoding and
high-mask quotient argument require it. Those arguments now apply
unchanged, proving canonical code recovery and membership.

Conversely, take the positive witnesses constructed for the
109-operation system. They have $e<q$, $\lambda>q$, and $S_3>0$.
Set $\Omega=Z\lambda-2e>0$ and replace the slack by
$\alpha=q^2-Y>0$. All other witnesses remain valid.

Only the addition of $\mathcal C^2$ to $Y$ disappears.
The coefficient subtraction, its product with $\mathcal C^2$, and
the subtraction defining $S_3$ have the same costs as before.
The positive unknown $\Omega$ introduces one free equality and no
arithmetic instruction beyond these existing registers.
The complete check is \file{round10\_1980\_certificate.py}.
\end{proof}

\subsection{Factoring the odd Pell root}

For this independent change to the 109-operation system, write
\[
\begin{gathered}
 N=n,\quad R=r,\quad U=wN^2,\quad Y_P=sN^2,
 \quad M=RY_P,\\
 a=M(U+1),\quad Q=UM^2,
 \quad P=2Q+1,\quad J=2R+1.
\end{gathered}
\]
The quantity $Q$ is already computed. The old first Pell norm
and index equations were
\[
 \bigl((2Q+1)^2-1\bigr)k^2+1=\tau_{\rm old}^2,
 \qquad k=R+1+2Qh_{\rm old}.
\]
Replace them by
\begin{equation}\label{eq:triangular-pell}
 \tau(\tau+1)=Q(Q+1)k^2,
 \qquad k=R+1+hQ,
 \qquad \tau,h>0.
\end{equation}
The mathematical Pell parameter $P=2Q+1$ need not be computed
by the new certificate.

\begin{proposition}\label{prop:triangular-pell108}
Equations \eqref{eq:triangular-pell} give a positive-domain
equivalent system with 108 arithmetic instructions:
59 multiplications and 49 additions. This independent variant
retains 33 positive unknowns and 21 equations.
\end{proposition}
\begin{proof}
For necessity, $P$ is odd, and the recurrence for $\chi_P(t)$
shows that every such coordinate is odd. The old positive norm
has $\tau_{\rm old}>1$. Thus
\[
 \tau=(\tau_{\rm old}-1)/2\ge1,
 \qquad h=2h_{\rm old}>0
\]
are integers satisfying \eqref{eq:triangular-pell}.
The factored norm is simply the old norm identity divided by four.
The certificate verifies the new polynomial equation directly;
this witness map does not introduce an uncharged division instruction.

For sufficiency, the initial coding bounds give $R\ge N\ge8$,
so $U,Y_P\ge64$, $M\ge512$, and $Q>R+1$.
Put $\widehat\tau=2\tau+1$. The new norm implies
\[
 \widehat\tau^2-(P^2-1)k^2=1,
 \qquad \widehat\tau>1.
\]
The new congruence gives $k\ge R+2$. The unchanged positive
interval gives $c>Y_Pk\ge64(R+2)>J$. The signed Pell block
therefore gives $c=\psi_a(J)$ and $d=\chi_a(J)$.

There is a positive index $t$ with
$k=\psi_P(t)$ and $\widehat\tau=\chi_P(t)$.
The elementary congruence $\psi_P(t)\equiv t\pmod{P-1}$
also holds modulo $Q$. Since $0<R+1<Q$, the new index equation
implies
\[
 t=R+1+vQ,\qquad v\ge0.
\]
If $v\ge1$, the elementary Pell growth estimates give
\[
 \frac ck\le
 \frac{(2a)^{2R}}{(2P-1)^{R+Q}}
 \le\left(\frac{2a}{2P-1}\right)^{2R}<\frac12.
\]
Here $Q\ge R$, and
$2a/(2P-1)=2M(U+1)/(4UM^2+1)<1/2$ follows from
$U\ge64$ and $M\ge512$. This contradicts $c/k>Y_P\ge64$.
Therefore $v=0$ and $k=\psi_P(R+1)$.

Applying the same Pell congruence now modulo $2Q$ gives
$k\equiv R+1\pmod{2Q}$. Since $k-R-1=hQ$, the new positive
integer $h$ is even. Consequently
\[
 h_{\rm old}=h/2>0,\qquad
 \tau_{\rm old}=2\tau+1>0
\]
restore both old equations, preserving every other witness.
This argument precedes the exponent relations and binomial rounding;
it does not assume either of them to prove the required evenness.

After $Q$ is available, the old norm block used seven operations:
form $2Q$, add two, multiply these two values, square $k$,
multiply, add one, and square $\tau_{\rm old}$.
The new block uses six: form $Q+1$ and $Q(Q+1)$, square $k$,
multiply, form $\tau+1$, and multiply $\tau(\tau+1)$.
The index product $hQ$ has the same cost as the former $2Qh_{\rm old}$,
so $2Q$ is no longer needed anywhere in the schedule.
The exact saving is one multiplication, as checked by
\file{round11\_1980\_pell\_certificate.py}.
\end{proof}

\subsection{Combining the two savings}

\begin{theorem}\label{thm:best107}
For each r.e.\ set, the fixed admissible input-unit encoding gives
a membership-equivalent system in 34 positive unknowns and 22 equations
whose straight-line certificate has \textbf{107 arithmetic instructions}:
59 multiplications and 48 additions. Equality tests and supplied
parameters are free. Generating the fixed numerals from $1$ adds
seven instructions, giving 114 operations.
\end{theorem}
\begin{proof}
Apply both replacements to the 109-operation system. Their
necessity maps affect disjoint witnesses: the coefficient-bound
change replaces $\alpha$ and supplies $\Omega$, while the Pell
change replaces $\tau$ and $h$. Thus both maps preserve the
other change's equations and positivity.

For soundness, first use $\Omega,\sigma>0$ to establish the
preliminary bounds $0<S<n$, $0\le T<n$, and $R\ge N\ge8$,
as in Proposition~\ref{prop:positive-coefficient108}.
The index proof in Proposition~\ref{prop:triangular-pell108}
then restores the old positive Pell witnesses without any exponent
or decoding assumption. The positive-coefficient proof can now
continue with those witnesses and recover membership. This order
establishes the equivalence without a circular appeal between the
two changes.

The combined schedule removes one addition and one multiplication
from 109 instructions. The coefficient change supplies the sole
extra positive unknown and free equality; the Pell change alters
neither count. The complete primitive list and all residual
identities are checked by \file{round12\_1980\_\allowbreak certificate.py}.
The numeral set and exponent $L=5^{16}$ are unchanged.
\end{proof}
